\documentclass[a5paper, 10pt]{article}


\usepackage{cite, amsmath, amssymb, booktabs, lastpage, graphicx, float}
\usepackage[hidelinks]{hyperref}

\usepackage{geometry}
\geometry{margin=1.5cm,a5paper,twoside,inner=2cm}

\usepackage{setspace}


\usepackage{amsthm}

%\usepackage[
%			type={CC},
%			modifier={by},
%			version={4.0},
%			imagewidth=5em,
%			imagedistance=5em,
%			]{doclicense}
%
%\makeatletter
%	\def\blfootnote{\gdef\@thefnmark{}\@footnotetext}
%\makeatother

\theoremstyle{plain}
\newtheorem{lemma}{Lemma}
\newtheorem{theorem}{Theorem}
\newtheorem{conjecture}{Conjecture}
\newtheorem{proposition}[theorem]{Proposition}
\newtheorem*{corollary}{Corollary}
\renewcommand{\qedsymbol}{\rule{.7em}{.7em}}

\theoremstyle{remark}
\newtheorem*{remark}{Remark}
\newtheorem*{note}{Note}
\newtheorem{case}{Case}

\theoremstyle{definition}
\newtheorem{definition}{Definition}
\newtheorem{example}{Example}

%% Changing the title, author, etc.
\usepackage{graphicx}
\makeatletter
\renewcommand{\maketitle}{
	\begin{center}
    \rule[.2em]{\textwidth}{0.353mm}
		\begin{minipage}[m]{0.35\textwidth}
			{\scriptsize
				\begin{center}
					\textsf{\textbf{\huge MATCH}\\
						\textit{Communications in Mathematical\\
							and in Computer Chemistry}
					}
			\end{center}}
		\end{minipage}\hfill
		\begin{minipage}[m]{0.65\textwidth}
			\begin{flushright}
				\baselineskip=10px
				{\scriptsize\sffamily{\itshape  MATCH Commun. Math. Comput. Chem.}
				\textbf{\vol{}}
				(\pubyear)
					\thepage--\pageref{LastPage}}\\
				{\scriptsize\sffamily {\bfseries ISSN:} 0340--6253}\\
				{\scriptsize\sffamily \textbf{doi:} \doi{10.46793/match.\vol-\num.\id}}
			\end{flushright}
		\end{minipage}
		\rule[1em]{\textwidth}{.353mm}
		\baselineskip=0.30in
		{\Large\bfseries \@title} \par
		\vspace{5mm}
		\baselineskip=0.2in
		{\large\bfseries \@author}\par
		\vspace{1mm}
		{\it \@address} \par
		{\small\tt \@email} \par
		\vspace{3mm}
		{\small (Received \@date)} \par
	\end{center}
	\vspace{3mm}
}
\newcommand{\address}[1]{\def\@address{#1}}
\newcommand{\email}[1]{\def\@email{#1}}
\newcommand{\doi}[1]{\href{https://doi.org/#1}{#1}}
%\title{}
%\author{}
%\address{}
%\email{}


\newcommand{\vol}{00}
\newcommand{\num}{00}
\newcommand{\pubyear}{0000}
\newcommand{\id}{xxxxx}
%%setting page-number
\setcounter{page}{1}


\thispagestyle{empty}
\makeatother

%% Making << \acknowledgment >> command
\newcommand{\acknowledgment}[1]{\vspace{5mm}\singlespacing
	{\noindent\textbf{\textit{Acknowledgment\/}:} #1}
}

%% Setting page numbers
\usepackage{fancyhdr}

%% now reset the headers and footers
\fancyhead{}
\fancyfoot{}
\renewcommand{\headrulewidth}{.3pt}
\voffset = 18pt
\headsep = 3pt

\renewcommand*{\thefootnote}{\fnsymbol{footnote}}

\usepackage[margin=1cm,%
			font=footnotesize,%
			format=hang,%
			labelsep=period,%
			labelfont=bf]{caption}

%%%%%%%%%%%%%%%%%%%%%%%%%%%%%%%%%%%%%%%%%%%%%%%%%%
%%%%%%%%%%%%%%%%%%%%%%%%%%%%%%%%%%%%%%%%%%%%%%%%%%
% PLACE ADDITIONAL PACKAGES HERE (if it is needed):

% \usepackage{package}
% \usepackage[options]{package}
\usepackage[T1]{fontenc}
\usepackage{lmodern}
\usepackage{mathtools}
\usepackage{array}
\usepackage{enumitem}
\setlist[itemize]{topsep=2pt,itemsep=2pt}
\setlist[enumerate]{topsep=2pt,itemsep=2pt}
\emergencystretch=3em
\hbadness=2000
\setlength{\marginparwidth}{2cm}

%%%%%%%%%%%%%%%%%%%%%%%%%%%%%%%%%%%%%%%%%%%%%%%%%%
%%%%%%%%%%%%%%%%%%%%%%%%%%%%%%%%%%%%%%%%%%%%%%%%%%

%% make the odd pages have the section name on the top right
\fancyhead[RO]{\footnotesize \thepage}
%% make the even pages have the chapter name on the top left
\fancyhead[LE]{\footnotesize \thepage}

%% making page-numbers visible
\pagestyle{fancy}




% YOU MAY ADD SOME ADDITIONAL LATEX PACKAGES TO THE PREAMBLE, BUT DO NOT MODIFY
% THE EXISTING CODE ABOVE THIS LINE.

% Custom macros for the Loyola index family
    \newcommand{\Wuzi}{LO(G;\alpha,\beta,\gamma)}
  \newcommand{\edgepsi}{\psi(x,y;\alpha,\beta,\gamma)}
  \newcommand{\bid}{\mathcal{F}_{\mathrm{BID}}}
% Path to figures
\graphicspath{{figures/}{./}{../figures/}}

% robust figure inclusion: render if present (root or figures/), else a small note
\newcommand{\safefig}[2]{%
  \IfFileExists{#1}{\includegraphics[width=#2\linewidth]{#1}}{%
   \IfFileExists{figures/#1}{\includegraphics[width=#2\linewidth]{#1}}{%
    \fbox{\parbox[c][3cm][c]{0.8\linewidth}{\centering Figure not in project. Upload \texttt{\detokenize{#1}}.}}}}%
}

%%%%%%%%%%%%%%%%%%%%%%%%%%%%%%%%%%%%%%%%%%%%%%%%%%%%%%%%%%%%%%%%%%%%%%%%%%%%%%%%%%%%%%%

\title{The Loyola Index Family and the Resolution Limit of Degree-Ratio Indices on Trees}
\author{Advik Natarajan$^{a}$, Ganesh Shiwakoti$^{b}$, Michael Arockiaraj$^{a,}$\footnote{Corresponding author.}}

\address{$^a$Department of Mathematics, Loyola College, Chennai, India\\
    $^b$Department of Mathematics and Statistics, Florida Atlantic University, USA
    }

\email{marockiaraj@gmail.com}

\date{\today}

% !TeX program = pdflatex

\begin{document}

\maketitle

\begin{abstract}
  Every vertex-degree-based index that sees an edge only through the
  ratio of its end-degrees, such as the geometric--arithmetic,
  symmetric division deg and degree-ratio Sombor indices, depends on a
  graph only through its histogram of normalised edge imbalances
  $q_{uv} = |d_u - d_v|/(d_u + d_v)$, a direct consequence of Rada's
  structure relation. By exhaustive enumeration we determine when this
  histogram loses degree-pair information: it determines the degree-pair
  counts of a tree up to order twelve, and of a chemical tree up to order
  fifteen, but not from order thirteen (sixteen for chemical trees), where
  explicit colliding trees exist at every order. These are the first
  orders at which all degree-ratio indices fail simultaneously; a single
  index, such as the symmetric division deg index, can fail earlier. As a
  test case we use the \emph{Loyola index family}
  \begin{equation*}
  LO(G;\alpha,\beta,\gamma)
  \;=\; \sum_{uv \in E(G)}
  (d_u d_v)^{\alpha}\,(d_u + d_v)^{\beta}\,
  e^{\gamma q_{uv}},
  \end{equation*}
  an exponential extension of the Gutman--Milovanovi\'c family. By Rada's
  Lindemann--Weierstrass theorem its pure-imbalance points at algebraic
  $\gamma \neq 0$ separate exactly the imbalance histograms, and two
  chemical graphs exchange order at most five times along $\gamma$.
  Pre-registered tests on measured boiling points did not demonstrate a
  predictive gain from $\gamma$; on those data ridge regression on the
  degree-pair counts had a higher cross-validated $Q^2$ than every single
  index evaluated, in point estimate. In an exploratory comparison, a gain
  on vaporisation enthalpies seen in a narrow search box shrinks to about
  zero once the parent is searched over a box that does not truncate it.
\end{abstract}

\onehalfspacing

\section{Introduction}\label{sec:intro}

  Degree-based topological indices run from the Randi\'c and Zagreb
  indices~\cite{Randic1975, GutmanTrinajstic1972} and their parametric
  extensions~\cite{BollobasErdos1998, ZhouTrinajstic2010} to the Sombor
  index~\cite{Gutman2021Sombor} and its many
  variants~\cite{Gutman2026SomborSurvey}. Extremal results on
  bond-incident-degree (BID) indices are surveyed
  in~\cite{AliGutmanFurtula2026BID}. On graphs with $\Delta \le 4$ every
  such index is a linear functional of the ten edge-degree-pair
  counts~\cite{Rada2026LinearVDB}, so no new index can carry more graph
  information than those counts. The question for a new index is
  therefore which count vectors it separates, and whether it predicts
  better than the counts themselves.

  A growing group of indices sees each edge only through the ratio of
  its end-degrees: the geometric--arithmetic and arithmetic--geometric
  indices, the symmetric division deg index~\cite{VukicevicGasperov2010},
  the diminished Sombor index~\cite{MovahediGutmanRedzepovicFurtula2026}
  and the degree-ratio
  indices~\cite{NyauliBuragohain2027DRSO, GutmanBuragohainNyauli2026},
  including the degree-ratio Sombor
  index~\cite{GutmanRedzepovicBuragohainNyauli2027, AliGutmanEtAl2027DRSO, ShaoJahanbani2027}.
  We show that all of them depend on a graph only through its histogram
  of normalised imbalances $q_{uv} = |d_u - d_v|/(d_u + d_v)$, and we
  determine when this histogram loses degree-pair information: for trees
  from order $13$, and from order $16$ for chemical trees.

  We study this limit through the three-parameter extension
  \[
  LO(G;\alpha,\beta,\gamma) = \sum_{uv \in E(G)}
  (d_u d_v)^{\alpha}(d_u+d_v)^{\beta} e^{\gamma q_{uv}}
  \]
  of the Gutman--Milovanovi\'c family $M_{\alpha,\beta}$
  \cite{GutmanMilovanovic2020AKCE, MolinaSigarreta2025GM}, which we call
  the \emph{Loyola index family}. It is a subfamily of the exponential
  vertex-degree-based indices of Rada~\cite{Rada2019ExpVDB}, and its
  pure-$\gamma$ points $LO(G;0,0,\gamma) = \sum_{uv} e^{\gamma q_{uv}}$
  are the exponential indices of the edge function $\gamma q_{uv}$. As
  $\gamma \to 0$, $LO(G;0,0,\gamma) = m + \tfrac{\gamma}{2}\,IRLA(G) + O(\gamma^2)$,
  where $IRLA(G) = 2\sum_{uv} q_{uv}$~\cite{RetiSharafdini2018IRLA} is a
  normalised form of Albertson's irregularity~\cite{Albertson1997, LinZhouWangMiao2022}. Its extremal theory is developed
  in~\cite{CruzMonsalveRada2020, Sigarreta2022ExpVDB, GaoAMC2024, HuMATCH2022},
  and parameter-optimised versions of the parent have already been fitted
  to octane properties~\cite{AliGutmanFurtulaAlbalahiHamza2025}. The
  variable Euler--Sombor index~\cite{GutmanRedzepovicKizilirmakKulli2025EU}
  and the cosine-rule Sombor indices~\cite{NyauliBuragohain2026CoRSO}
  weight the imbalance polynomially and without normalisation. The family
  is not proposed as a better predictor. Gutman's survey~\cite{Gutman2026SomborSurvey}
  observes that new variants hardly ever improve on earlier ones in
  applications, and our tests agree.

  Our contributions are as follows.
  \begin{enumerate}
  \item[(i)] The ratio-plane theorem (Theorem~\ref{thm:ratioplane}):
  on the plane $\beta = -2\alpha$ the index depends only on the
  $q$-histogram, the pure-$\gamma$ points at algebraic $\gamma \neq 0$
  separate exactly the $q$-histograms, and distinct degree-pair profiles
  with equal histograms occur among trees of every order from $13$
  (from $16$ for $\Delta \le 4$). That these orders are the first is
  established by exhaustive enumeration. The corollary extends the
  collisions to every degree-ratio index.
  \item[(ii)] At rational parent exponents a nonzero algebraic $\gamma$
  never destroys a distinction made by the parent
  (Proposition~\ref{prop:offplane}).
  \item[(iii)] A bound on the order exchanges along $\gamma$
  (Proposition~\ref{prop:laguerre}) and per-graph identifiability
  (Proposition~\ref{prop:identrank}). The extremal trees of the
  pure-$\gamma$ points (Proposition~\ref{prop:extremal}) follow from
  known criteria and are included with a short proof.
  \item[(iv)] Freeing $\gamma$ gives no demonstrated gain once the
  parent's search box does not truncate it, and on measured boiling points
  ridge regression on the degree-pair counts leads every single index
  evaluated in point estimate (Section~\ref{sec:sensitivity}). On the octanes it does not.
  \end{enumerate}

\section{Preliminaries}\label{sec:prelim}

Let $G = (V, E)$ be a finite simple connected graph with at least one
edge, of order $n$ and size $m$, with degrees $d_v$, maximum degree
$\Delta$ and minimum degree $\delta$. The chemometric section uses
hydrogen-suppressed alkane skeletons, for which $\Delta \le 4$; the
mathematical results do not assume this bound unless stated. For
integers $1 \le i \le j$ let $m_{ij}(G)$ be the number of edges with
unordered degree pair $(i,j)$; the vector $(m_{ij}(G))_{i \le j}$ is the
\emph{degree-pair profile} of $G$.

\begin{definition}[Loyola index family]\label{def:wuzi}
For $\alpha, \beta, \gamma \in \mathbb{R}$,
\begin{equation}\label{eq:wuzi}
\Wuzi
\;=\; \sum_{uv \in E(G)}
(d_u d_v)^{\alpha}\,(d_u + d_v)^{\beta}\,
\exp\!\left(\gamma\,\frac{|d_u - d_v|}{d_u + d_v}\right).
\end{equation}
\end{definition}

Write $q_{uv} = |d_u - d_v|/(d_u + d_v)$ and $q_{ij} = (j-i)/(i+j)$ for
$i \le j$. Each summand depends only on the degree pair, so
\begin{equation}\label{eq:pairsum}
LO(G;\alpha,\beta,\gamma) = \sum_{i \le j} m_{ij}(G)\,
(ij)^{\alpha}(i+j)^{\beta} e^{\gamma q_{ij}} .
\end{equation}
In particular $LO(G;\alpha,\beta,0) = M_{\alpha,\beta}(G)$, and $\gamma$
has no effect on regular graphs. Special cases, catalogued
in~\cite{MolinaSigarreta2025GM}, include $M_1 = LO(0,1,0)$,
$M_2 = LO(1,0,0)$, $HM = LO(0,2,0)$~\cite{ShirdelRezapourSayadi2013HM},
${}^{m}\!M_2 = LO(-1,0,0)$, the Randi\'c index $R = LO(-\tfrac12,0,0)$,
the sum-connectivity index
$\chi = LO(0,-\tfrac12,0)$~\cite{ZhouTrinajstic2009sumconn},
$H/2 = LO(0,-1,0)$, $ISI = LO(1,-1,0)$,
$GA = 2\,LO(\tfrac12,-1,0)$~\cite{VukicevicFurtula2009GA} and
$AG = \tfrac12 LO(-\tfrac12,1,0)$.

\begin{definition}[$q$-histogram]\label{def:qhist}
The $q$-histogram of $G$ is
$n_q(G) = \sum_{i \le j,\; q_{ij} = q} m_{ij}(G)$, $q \in [0,1)$.
\end{definition}

Equal profiles give equal $q$-histograms; the converse fails
(Theorem~\ref{thm:ratioplane}(iii)). Since $4ij/(i+j)^2 = 1 - q_{ij}^2$,
\begin{equation}\label{eq:factor}
LO(G;\alpha,\beta,\gamma) = 4^{-\alpha} \sum_{uv \in E(G)}
(d_u + d_v)^{2\alpha+\beta}\,(1 - q_{uv}^2)^{\alpha}\, e^{\gamma q_{uv}},
\end{equation}
so $\kappa = 2\alpha + \beta$ acts only on the degree sum, and $\alpha$,
$\gamma$ act only on $q_{uv}$, a function of the degree ratio.

\begin{definition}[Ratio plane]\label{def:ratioplane}
The ratio plane is
$\Pi = \{(\alpha,\beta,\gamma) : \beta = -2\alpha\}$. It contains the
pure-$\gamma$ points $(0,0,\gamma)$, $GA$ and $AG$.
\end{definition}

\section{The resolution limit}\label{sec:limit}

\subsection{The degree-pair ceiling}\label{sec:ceiling}

No vertex-degree-based index separates graphs with equal profiles, so the
number of distinct profiles is a ceiling on distinct values.

  \begin{proposition}[Degree-pair ceiling]\label{prop:collision}
  If $m_{ij}(G) = m_{ij}(H)$ for all $i \le j$, then
  $LO(G;\theta) = LO(H;\theta)$ for every $\theta \in \mathbb{R}^3$, and
  $TI(G) = TI(H)$ for every vertex-degree-based index $TI$. Among the
  $18$ octane isomers exactly $16$ distinct count vectors occur, and the
  pairs $\{$3-methylheptane, 4-methylheptane$\}$ and
  $\{$3,4-dimethylhexane, 3-ethyl-2-methylpentane$\}$ share their
  count vectors but differ in measured properties
  (Table~\ref{tab:collisions}). Hence every vertex-degree-based
  descriptor, tuned or not, carries a nonzero irreducible in-sample error on this benchmark.
  \end{proposition}

  \begin{proof}
  By~\eqref{eq:pairsum}, $LO(G;\theta)$ depends on $G$ only through
  $(m_{ij}(G))_{i \le j}$, and so does every vertex-degree-based
  index~\cite{Rada2026LinearVDB}, since these graphs share the
  M-polynomial~\cite{DeutschKlavzar2015}. The octane counts are computed
  and cross-checked in the reproducibility package~\cite{LoyolaRepo}.
  \end{proof}

For the octane boiling points the resulting floor on the sum of squared
errors of any fixed function of the counts is $2.925\,^{\circ}$C$^2$, a
root mean square of $0.403\,^{\circ}$C. It bounds a single fixed
predictor, not the cross-validated errors of Section~\ref{sec:sensitivity},
where each fold fits its own.

\begin{table}[htbp]
\centering
\caption{The two degree-pair collision pairs among the $18$ octane
isomers, indistinguishable by every vertex-degree-based index. Units:
$T_B$ in $^{\circ}$C, $\Delta H_f$ and $\Delta H_{\mathrm{vap}}$ in
kcal\,mol$^{-1}$, $S$ in cal\,mol$^{-1}$\,K$^{-1}$; $\omega$ is
dimensionless.}
\label{tab:collisions}
\small
\begin{tabular}{lrrrrr}
\toprule
Isomer & $T_B$ & $\Delta H_f$ & $\Delta H_{\mathrm{vap}}$ & $S$ & $\omega$\\
\midrule
\multicolumn{6}{l}{\emph{shared count vector} $2(1,2),\,1(1,3),\,2(2,2),\,2(2,3)$:}\\
3-methylheptane & $118.9$ & $-50.82$ & $9.52$ & $111.26$ & $0.371$\\
4-methylheptane & $117.7$ & $-50.69$ & $9.48$ & $109.32$ & $0.372$\\
\midrule
\multicolumn{6}{l}{\emph{shared count vector} $2(1,2),\,2(1,3),\,2(2,3),\,1(3,3)$:}\\
3,4-dimethylhexane & $117.7$ & $-50.91$ & $9.32$ & $106.59$ & $0.340$\\
3-ethyl-2-methylpentane & $115.6$ & $-50.48$ & $9.21$ & $106.06$ & $0.330$\\
\bottomrule
\end{tabular}
\end{table}

The ceiling is attained at almost every parameter point.

  \begin{proposition}[Generic discrimination]\label{prop:generic}
  Let $m \neq m'$ be two profiles supported on the pairs
  $P_{\Delta} = \{(i,j) : 1 \le i \le j \le \Delta\}$, and write
  $z_{ij} = (\ln(ij), \ln(i+j), q_{ij})$. Then
  $D(\theta) = \sum_{(i,j) \in P_{\Delta}} (m_{ij} - m'_{ij})\,
  e^{\theta^{\top} z_{ij}}$ is a nonzero real-analytic function, whose zero
  set has Lebesgue measure zero in $\mathbb{R}^3$~\cite{Mityagin2020}.
  Hence graphs with distinct profiles are separated by
  $LO(\cdot\,;\theta)$ for almost every $\theta$.
  \end{proposition}

  \begin{proof}
  The map $(i,j) \mapsto z_{ij}$ is injective, since $ij$ and $i+j$
  determine the pair. Choose $v$ with the $v^{\top} z_{ij}$ pairwise
  distinct, which excludes only finitely many hyperplanes of $v$. Along $\theta = tv$, $D$ is an exponential sum with distinct
  exponents and coefficients not all zero, so it is not identically zero.
  \end{proof}

  The proof uses only the injectivity of
  $(i,j) \mapsto (\ln(ij), \ln(i+j))$, so it applies to the parent on the
  $(\alpha,\beta)$-plane: generic discrimination is inherited, not conferred
  by $\gamma$. As there are countably many pairs of profiles, almost every
  point of the family is weakly discriminating, in the sense of Rada,
  Rodr\'iguez and Sigarreta~\cite{RadaRodriguezSigarreta2022Weak}, over
  every class at once.

\subsection{The ratio plane}\label{sec:ratioplane}

  Rada~\cite{Rada2026LinearVDB} shows that no vertex-degree-based index
  separates trees with equal profiles. The theorem below concerns a
  coarser loss: on $\Pi$, distinct profiles with equal $q$-histograms
  receive equal values.

  \begin{theorem}[Ratio plane]\label{thm:ratioplane}
  Let $G$ and $H$ be graphs.
  \begin{enumerate}
  \item[(i)] On $\Pi$,
  $LO(G;\alpha,-2\alpha,\gamma) = \sum_q n_q(G)\,\bigl((1-q^2)/4\bigr)^{\alpha}
  e^{\gamma q}$ depends on $G$ only through its $q$-histogram, and $G$ and
  $H$ agree at every point of $\Pi$ if and only if $n_q(G) = n_q(H)$
  for all $q$. Otherwise their values differ at every point of $\Pi$
  outside a set of measure zero.
  \item[(ii)] For every nonzero algebraic $\gamma$,
  $LO(G;0,0,\gamma) = LO(H;0,0,\gamma)$ if and only if
  $n_q(G) = n_q(H)$ for all $q$. On any finite set of graphs the same
  holds for all but finitely many real $\gamma$.
  \item[(iii)] Let $G$ and $H$ be trees with distinct
  edge-degree-pair profiles and equal $q$-histograms, each having a
  leaf adjacent to a vertex of the same degree $d \ge 2$. Attaching a pendant path of length $L \ge 1$ at such a leaf in each
  produces trees with the same two properties and the same maximum
  degree. The profiles change by the same vector, whether or not the
  chosen leaf edge was the only $(1,d)$-edge.
  Consequently, for every $n \ge 13$ (and every $n \ge 16$ among trees
  with $\Delta \le 4$) there are non-isomorphic trees of order $n$
  with distinct degree-pair profiles that no point of $\Pi$ separates.
  \end{enumerate}
  \end{theorem}

  \begin{proof}
  (i) The formula follows from~\eqref{eq:factor} with $\beta = -2\alpha$.
  If the histograms differ, the difference at $\alpha = 0$ is
  $\sum_q (n_q(G) - n_q(H))\,e^{\gamma q}$, an exponential sum with
  distinct exponents and not all coefficients zero, which is not
  identically zero because the term with the largest exponent dominates
  as $\gamma \to +\infty$. The difference on $\Pi$ is therefore a nonzero
  real-analytic function of $(\alpha,\gamma)$, whose zero set has
  measure zero~\cite{Mityagin2020}.
  (ii) Every imbalance is rational, so for algebraic $\gamma \neq 0$ the
  exponents $\gamma q$ are distinct algebraic numbers. By the
  Lindemann--Weierstrass theorem~\cite{Baker1975} the $e^{\gamma q}$ are
  linearly independent over the algebraic numbers, so the integer
  coefficients $n_q(G) - n_q(H)$ vanish. For real $\gamma$, a nonzero
  exponential sum with $k$ distinct exponents has at most $k - 1$ real
  zeros (Proposition~\ref{prop:laguerre}).
  (iii) The attachment, a device of Rada~\cite{Rada2026LinearVDB},
  changes each profile by $-(1,d) + (2,d) + (L-1)(2,2) + (1,2)$, so
  distinct profiles stay distinct and equal histograms stay equal; the
  new vertices have degree at most $2$. It remains to give the starting
  pairs. At order $13$: $T_1$ is a path $c\,a\,b\,d\,e\,\ell$ in which $c$
  carries five pendant leaves, $b$ and $d$ one each and $\ell$ is a leaf,
  with edges $(1,6)^{5}, (2,6), (2,3)^{2}, (3,3), (1,3)^{2}, (1,2)$;
  $T_2$ is a path $c\,b\,x\,y\,z$ in which $c$ carries five pendant
  leaves, $b$ one and $z$ two, with edges
  $(1,6)^{5}, (3,6), (1,3)^{3}, (2,3)^{2}, (2,2)$. Both have the
  $q$-histogram $\{0^{1}, \tfrac15^{2}, \tfrac13^{1}, \tfrac12^{3}, \tfrac57^{5}\}$
  and a leaf adjacent to a vertex of degree $3$. At order $16$ with
  $\Delta \le 4$: the profile
  $\{(1,2)^{2}, (1,4)^{8}, (2,4)^{2}, (4,4)^{3}\}$ is realised by a path
  $A\,B\,C\,D$ of vertices of degree $4$, where $A$ and $D$ carry three
  pendant leaves each and $B$ and $C$ one pendant leaf and one pendant
  path of length $2$ each; the profile
  $\{(1,4)^{8}, (2,2)^{3}, (2,4)^{4}\}$ is realised by vertices $A, B, C$
  of degree $4$ joined as $A\,w\,B$ and $B\,v_1 v_2 v_3 v_4\,C$ through
  vertices of degree $2$, with three, two and three pendant leaves. Both
  have the $q$-histogram $\{0^{3}, \tfrac13^{4}, \tfrac35^{8}\}$ and a
  leaf adjacent to a vertex of degree $4$.
  \end{proof}

  Part~(ii) is Theorem~3.4 of~\cite{Rada2019ExpVDB} applied to the
  algebraic edge weights $\gamma q_{ij}$, whose classes of equal weight are
  the $q$-classes. The construction in part~(iii) parallels Example~2.2
  of~\cite{RadaRodriguezSigarreta2022Weak}, where indices that depend
  only on the degree sum are defeated by chemical trees whose counts agree
  on classes of equal sum; here the classes are those of equal ratio, which
  for $\Delta \le 4$ are $\{(1,2),(2,4)\}$ and $\{(2,2),(3,3),(4,4)\}$ (ignoring
  $(1,1)$, which occurs only in $K_2$).
  Song and Huang~\cite{SongHuang2023EDV} determined the trees of order at
  most $10$ with equal edge division vectors, which receive equal values of
  every edge-additive eccentric index; neither that partition of trees nor
  ours refines the other. We found no earlier determination of the orders $13$ and
  $16$, or of explicit colliding trees for degree-ratio indices, in a
  literature search recorded in the repository~\cite{LoyolaRepo}.

  \begin{corollary}[Degree-ratio indices]\label{cor:ratioclass}
  Let $\rho_{uv} = \min(d_u,d_v)/\max(d_u,d_v)$ and
  $I_f(G) = \sum_{uv \in E(G)} f(\rho_{uv})$ for any function $f$ on
  $(0,1]$. Since $\rho_{uv} = (1-q_{uv})/(1+q_{uv})$, $I_f$ depends on $G$
  only through its $q$-histogram, and the trees of
  Theorem~\ref{thm:ratioplane}(iii) are not separated by any such index.
  This applies to $GA$, $AG$, the exponential geometric--arithmetic
  index~\cite{DasMondal2024ExpGA}, the diminished Sombor
  index~\cite{MovahediGutmanRedzepovicFurtula2026}, the symmetric division
  deg index~\cite{VukicevicGasperov2010}, the inverse symmetric division
  deg index and its exponential~\cite{DasMondalRazaPal2025ExpISDD}, and
  the degree-ratio
  indices~\cite{NyauliBuragohain2027DRSO, GutmanBuragohainNyauli2026},
  including the degree-ratio Sombor
  index~\cite{GutmanRedzepovicBuragohainNyauli2027, AliGutmanEtAl2027DRSO, ShaoJahanbani2027}.
  \end{corollary}

  The corollary is the structure observation of~\cite[Proposition~2.5]{Rada2019ExpVDB}
  for edge weights that depend only on the degree ratio; what is new are
  the explicit witnesses and the orders at which they first occur.
  Within this class the pure-$\gamma$ points at algebraic $\gamma \neq 0$
  attain the most possible, by part~(ii). $GA$ does not
  (Table~\ref{tab:firstfail}). The parent attains it as well: at
  $\gamma = 0$ part~(i) gives an exponential sum in $\alpha$ with distinct
  bases, which separates two distinct histograms at all but finitely many
  $\alpha$. The limit thus belongs to the dependence on the degree ratio,
  not to $\gamma$. The orders $13$ and $16$ are the first at which the $q$-histogram, and
  hence every degree-ratio index simultaneously, fails to separate
  distinct profiles. A single index can fail earlier through coincidences
  among its values. The symmetric division deg index merges two distinct
  profiles of chemical trees already at order $9$, for instance
  $\{(1,2), (1,3)^{2}, (2,2)^{4}, (2,3)\}$ and
  $\{(1,2)^{3}, (1,3), (2,3)^{3}, (3,3)\}$, both with value $58/3$, and the
  inverse symmetric division deg index does so at order $10$ (exact
  enumeration in the repository).

\subsection{Where the limit binds}\label{sec:degeneracy}

By exhaustive enumeration of all trees of orders $2$--$17$ in exact
arithmetic (Table~\ref{tab:firstfail}), two trees of order at most $12$
with equal $q$-histograms have equal degree-pair profiles, and the same
holds for trees with $\Delta \le 4$ of order at most $15$. With
Theorem~\ref{thm:ratioplane}(iii), the pure-$\gamma$ points at algebraic
$\gamma \neq 0$ separate all distinct profiles exactly up to those
orders and at no larger order. Among all trees, at orders $13$--$17$,
generic points off $\Pi$ (Proposition~\ref{prop:generic}) and the Randi\'c
index separate more trees than any point of $\Pi$. Degeneracy~\cite{Konstantinova1996Degeneracy},
measured as $100(1 - k/N)$ with $k$ distinct values among $N$ trees, is
$25.5\%$ for the pure-$\gamma$ points on the $106$ trees of order $10$
($79$ distinct values, the number of distinct profiles and so the floor for
any degree-pair-based index), tied with $R$, $GA$
and $AG$, while $M_1$ ($83.0\%$) and $M_2$ ($67.0\%$) are the most
degenerate. $GA$ can do worse than the pure-$\gamma$ points: it merges two
distinct $q$-histograms among trees first at order $17$, for instance
$\{(1,5)^{10}, (2,2)^{2}, (5,5)^{2}, (2,5), (1,2)\}$ and
$\{(4,5)^{3}, (1,4)^{5}, (1,5)^{6}, (2,5), (1,2)\}$.

\begin{table}[htbp]
\centering
\caption{First tree order at which discrimination is lost, from all trees
of orders $2$ to $17$. ``Below floor'': first order with fewer distinct
values than distinct profiles (values/profiles). ``Merges'': first order
at which an index on $\Pi$ gives equal values to two trees with
different $q$-histograms. Every count, ties included, is certified
exactly: for the pure-$\gamma$ points by Theorem~\ref{thm:ratioplane}(ii),
and for $R$ and $GA$ by exact arithmetic in the radicals
$(ij)^{a}(i+j)^{b}$~\cite{Besicovitch1940}.}
\label{tab:firstfail}
\small
\setlength{\tabcolsep}{3pt}
\begin{tabular}{lcccc}
\toprule
 & \multicolumn{2}{c}{all trees} & \multicolumn{2}{c}{trees with $\Delta \le 4$}\\
\cmidrule(lr){2-3}\cmidrule(lr){4-5}
Index & below floor & merges & below floor & merges\\
\midrule
$LO(0,0,\gamma)$ & $13$ ($566/570$) & never$^{a}$ & $16$ ($809/810$) & never$^{a}$\\
$R$ & $14$ ($1095/1100$) & n/a$^{b}$ & $17$ ($1191/1194$) & n/a$^{b}$\\
$GA$ & $13$ ($566/570$) & $17$ & $16$ ($809/810$) & $> 17^{c}$\\
\bottomrule
\end{tabular}

\smallskip
\begin{minipage}{0.95\linewidth}\footnotesize
$^{a}$~For every nonzero algebraic $\gamma$
(Theorem~\ref{thm:ratioplane}(ii)). $^{b}$~$R$ is not on $\Pi$.
$^{c}$~No merge up to order $17$, the largest examined.
\end{minipage}
\end{table}

Off $\Pi$ the imbalance factor can restore discrimination lost at
particular rational exponents, and at rational exponents a nonzero
algebraic $\gamma$ never destroys any.

  \begin{proposition}[Parent distinctions survive]\label{prop:offplane}
  Let $\alpha, \beta \in \mathbb{Q}$ and let $\gamma \ne 0$ be algebraic.
  For graphs $G$, $H$ and each imbalance value $q$ write
  $c_q(G,H) = \sum_{(i,j):\, q_{ij} = q} \bigl(m_{ij}(G) - m_{ij}(H)\bigr)
  (ij)^{\alpha} (i+j)^{\beta}$. Then
  $LO(G;\alpha,\beta,\gamma) = LO(H;\alpha,\beta,\gamma)$ if and only if
  $c_q(G,H) = 0$ for every $q$. In particular, if
  $M_{\alpha,\beta}(G) \ne M_{\alpha,\beta}(H)$ then
  $LO(G;\alpha,\beta,\gamma) \ne LO(H;\alpha,\beta,\gamma)$.
  \end{proposition}

  \begin{proof}
  The difference is $\sum_q c_q(G,H)\, e^{\gamma q}$. For rational
  $\alpha, \beta$ the coefficients $c_q(G,H)$ are algebraic, and the
  exponents $\gamma q$ are distinct algebraic numbers, so by the
  Lindemann--Weierstrass theorem~\cite{Baker1975} the sum vanishes only
  if every $c_q(G,H)$ vanishes. Finally
  $\sum_q c_q(G,H) = M_{\alpha,\beta}(G) - M_{\alpha,\beta}(H)$.
  \end{proof}

  The algebraic hypothesis cannot be dropped: at $\gamma = \ln 4$,
  $LO(K_4;0,0,\gamma) = 6 = 3e^{\gamma/2} = LO(K_{1,3};0,0,\gamma)$. On
  all trees of order $16$ ($4069$ profiles), $LO(0,\tfrac12,0)$ and
  $LO(0,\tfrac13,0)$ take only $3475$ and $3476$ distinct values, while
  $LO(0,\tfrac12,1)$ and $LO(0,\tfrac13,1)$ separate all $4069$. These counts are exact (at
  $\gamma = 1$ by Proposition~\ref{prop:offplane}), but the effect is checked
  only up to order $16$ and not proved for larger orders, and it
  yields no predictive gain over a tuned parent
  (Section~\ref{sec:sensitivity}). The
  sum-connectivity index $\chi = M_{0,-1/2}$ is not weakly discriminating
  over chemical trees~\cite[Example~2.2]{RadaRodriguezSigarreta2022Weak};
  for the trees $S$, $T$ used there, at $(\alpha,\beta) = (0,-\tfrac12)$
  the difference $D(\gamma) = LO(S) - LO(T)$ satisfies $D(0) = 0$ and
  $D(1) \approx 0.2255$, but $D(10^{-6}) \approx 2.4 \times 10^{-7}$, so
  the separation is small for small $\gamma$.

\section{Behaviour along the imbalance axis}\label{sec:geometry}

  Write $\theta = (\alpha, \beta, \gamma)$ and
  $z_{uv} = \bigl(\ln(d_u d_v),\, \ln(d_u + d_v),\, q_{uv}\bigr)$, so that
  $LO_G(\theta) = \sum_{uv} e^{\theta^{\top} z_{uv}}$ and
  $\log LO_G$ is a log-sum-exp function of $\theta$.

  \begin{lemma}[Affine independence of the imbalance direction]\label{prop:indep}
  For $\Delta \ge 3$ the augmented vectors
  $\bigl(1, \ln(ij), \ln(i+j), q_{ij}\bigr)$, $(i,j) \in P_{\Delta}$,
  span $\mathbb{R}^4$. Consequently no choice of $(\alpha', \beta')$ and
  scaling constant reproduces the action of $\gamma \neq 0$ on all
  admissible degree pairs. For $\Delta \le 2$ it is reproduced exactly.
  \end{lemma}

  \begin{proof}
  On $(1,1), (2,2), (1,2)$ the vectors $(1, \ln(ij), \ln(i+j))$ have
  determinant $\ln 2 \cdot \ln(9/8) \ne 0$, so it suffices that $q_{ij}$
  is not affine in $(\ln(ij), \ln(i+j))$. If
  $q_{ij} = a + b \ln(ij) + c \ln(i+j)$ on $(1,1), (2,2), (1,2), (1,3)$,
  the first two pairs give $a = -c \ln 2$ and $b = -c/2$, and the last two
  give $\tfrac13 = c\,(\ln 3 - \tfrac32 \ln 2)$ and
  $\tfrac12 = c\,(\ln 2 - \tfrac12 \ln 3)$, so $c \ne 0$ and
  $8 \ln 3 = 13 \ln 2$, i.e.\ $3^{8} = 2^{13}$, which is false. Without
  the pair $(1,1)$, which occurs only in $K_2$, the vectors of
  $(1,2), (1,3), (2,2), (3,3)$ have determinant
  $-\tfrac16 \ln\tfrac32 \,\ln\bigl(2^{13}/3^{8}\bigr) \approx -0.0150 \ne 0$.
  If $e^{\theta^{\top} z} = c\, e^{\theta'^{\top} z}$ on $P_{\Delta}$, then
  $(\ln c, \theta' - \theta)$ is orthogonal to every augmented vector,
  hence zero, so $\theta = \theta'$. For $\Delta = 2$ the three vectors
  $(1, \ln(ij), \ln(i+j))$ of $(1,1), (1,2), (2,2)$ already have rank $3$,
  so $q_{ij}$ is an affine function of $\ln(ij)$ and $\ln(i+j)$ on them and
  $e^{\gamma q_{ij}}$ equals a constant times $(ij)^{a}(i+j)^{b}$; for
  $\Delta = 1$, $q \equiv 0$.
  \end{proof}

  \begin{theorem}[Log-convexity]\label{thm:logconvex}
  Let $F_G(\theta) = \log LO_G(\theta)$ and let $\mathbb{E}_{\theta}$,
  $\operatorname{Cov}_{\theta}$ refer to the edge probability measure
  $w_{uv}(\theta) = e^{\theta^{\top} z_{uv}} / LO_G(\theta)$. Then
  $\nabla F_G = \mathbb{E}_{\theta}[z]$ and
  $\nabla^2 F_G = \operatorname{Cov}_{\theta}(z) \succeq 0$, so $F_G$ is
  convex; it is strictly convex if and only if the $z_{uv}$ do not lie in
  a proper affine subspace of $\mathbb{R}^3$. In particular
  $\log LO$ is non-decreasing and convex in $\gamma$, and strictly
  increasing unless every edge joins vertices of equal degree.
  \end{theorem}

  \begin{proof}
  Differentiation gives the gradient and Hessian, and $w(\theta)$ has full
  support, so $\operatorname{Cov}_{\theta}(z)$ is singular in a direction
  $v$ if and only if $v^{\top} z_{uv}$ is constant over the edges. If no
  such $v$ exists the Hessian is positive definite and $F_G$ is strictly
  convex; if one exists, $t \mapsto F_G(\theta + tv)$ is affine and $F_G$
  is not strictly convex. Finally
  $\partial_{\gamma} F_G = \mathbb{E}_{\theta}[q] \ge 0$, with equality
  if and only if $q \equiv 0$.
  \end{proof}

  \begin{proposition}[Order exchanges along $\gamma$]\label{prop:laguerre}
  Fix $(\alpha,\beta)$ and graphs $G$, $H$, let
  $q_1 < \dots < q_r$ be the distinct imbalances of the degree pairs
  occurring in $G$ or $H$, and write
  $D(\gamma) = LO(G;\alpha,\beta,\gamma) - LO(H;\alpha,\beta,\gamma)
  = \sum_{k=1}^{r} c_k e^{\gamma q_k}$ with
  $c_k = \sum_{q_{ij} = q_k} (m_{ij}(G) - m_{ij}(H))(ij)^{\alpha}(i+j)^{\beta}$.
  Let $V$ be the number of sign changes in $(c_1,\dots,c_r)$ after
  zero terms are removed. If $D \not\equiv 0$, then $D$ has at most
  $V \le r - 1$ real zeros counted with multiplicity, and their number
  has the parity of $V$. In particular, for graphs with $\Delta \le 4$
  the possible imbalances are $0, \tfrac17, \tfrac15, \tfrac13,
  \tfrac12, \tfrac35$, and two such graphs exchange order at most five times as
  $\gamma$ ranges over $\mathbb{R}$, for every $(\alpha,\beta)$.
  \end{proposition}

  \begin{proof}
  With $x = e^{\gamma}$, $D$ is the generalised polynomial
  $\sum_k c_k x^{q_k}$, and Laguerre's extension of Descartes' rule of
  signs~\cite{Jameson2006} bounds its positive zeros by $V$, with the
  parity of $V$.
  \end{proof}

  The bound is attained. The trees of order $12$ with profiles
  $\{(1,3)^{2},\allowbreak (1,4)^{2},\allowbreak (1,6)^{5},\allowbreak (3,4),\allowbreak (4,6)\}$
  and
  $\{(1,2),\allowbreak (1,5)^{4},\allowbreak (1,6)^{4},\allowbreak (2,6),\allowbreak (5,6)\}$
  have $V = 5$, and at $(\alpha,\beta) = (-1,0)$ they cross at
  $\gamma \approx -17.2$, $-12.0$, $-2.0$, $2.7$ and $31.1$. Whether five
  crossings are attained for $\Delta \le 4$ is open.

  For a graph $G$ let $R(G)$ be the set of degree pairs it realises,
  $r(G)$ the number of distinct imbalances among them, and $A(G)$ the
  affine dimension of $\{z_{ij} : (i,j) \in R(G)\}$.

  \begin{proposition}[Per-graph identifiability]\label{prop:identrank}
  \begin{enumerate}
  \item[(i)] $A(G)$ equals the rank of $\operatorname{Cov}_{\theta}(z)$
  at every $\theta$, and the directions $v$ for which there is a constant
  $c$ with $LO(G;\theta + tv) = e^{ct}\,LO(G;\theta)$ for all $t$ and
  $\theta$ form a subspace of dimension $3 - A(G)$.
  \item[(ii)] $A(G) \le \min\bigl(|R(G)| - 1,\ r(G),\ 3\bigr)$.
  \item[(iii)] If $r(G) = 2$ with imbalances $q_1 \neq q_2$, then
  $v = (1, -2, c)$ with
  $c = [\ln(1 - q_2^2) - \ln(1 - q_1^2)]/(q_1 - q_2)$ is such a
  direction.
  \end{enumerate}
  \end{proposition}

  \begin{proof}
  (i) As in Theorem~\ref{thm:logconvex}, $\operatorname{Cov}_{\theta}(z)$
  is singular exactly in the directions $v$ with $v^{\top} z$ constant on
  $R(G)$, a subspace of dimension $3 - A(G)$, and these are exactly the
  rescaling directions, since $v^{\top}\nabla F_G = c$ is constant if
  and only if $\operatorname{Cov}_{\theta}(z)\,v = 0$ for all $\theta$.
  (ii) By~\eqref{eq:factor},
  $z_{ij} = (2s + \ln(1-q^2) - \ln 4,\ s,\ q)$ with $s = \ln(i+j)$, so
  pairs of equal imbalance lie on one line parallel to $(2,1,0)$, and $r$
  parallel lines lie in an affine space of dimension at most $r$.
  (iii) $v^{\top} z_{ij} = \ln(1 - q_{ij}^2) - \ln 4 + c\,q_{ij}$ takes
  the same value at $q_1$ and $q_2$.
  \end{proof}

  Equality holds in~(ii) for every set of admissible pairs with
  $\Delta \le 6$ (by exhaustive enumeration), hence for all molecular
  graphs. On the octanes the union of realised pairs contains nine of the
  ten admissible pairs and has affine dimension $3$, so the octane fits
  of Section~\ref{sec:sensitivity} are free of exact confounding of this
  kind. On the spider with four legs of length three, which realises
  only the imbalances $\tfrac13$ and $0$, $\gamma$ is confounded with
  $(\alpha,\beta)$ along the direction of part~(iii).

\begin{proposition}[Extremal trees of the pure-$\gamma$ points]\label{prop:extremal}
Let $n \ge 4$ and $\gamma > 0$. Among all trees of order $n$, the
path $P_n$ is the unique minimiser and the star $S_n$ the unique
maximiser of $LO(\cdot\,;0,0,\gamma)$, with
\[
(n-3) + 2e^{\gamma/3} \;\le\; LO(T;0,0,\gamma) \;\le\; (n-1)\,e^{\gamma(n-2)/n}.
\]
For $\gamma < 0$ the roles of $P_n$ and $S_n$ are exchanged.
\end{proposition}

\begin{proof}
Let $\gamma > 0$. Every edge of a tree of order $n$ has
$q_{uv} \le (n-2)/n$, with equality if and only if it is a
$(1,n-1)$-edge, which gives the upper bound and $S_n$. A tree with
$k \ge 2$ leaves has $k$ edges incident to a leaf, each with
$q \ge \tfrac13$, since the other end has degree $d \ge 2$ and
$(d-1)/(d+1)$ is increasing in $d$. Hence
$LO(T;0,0,\gamma) \ge (n-1-k) + k e^{\gamma/3} \ge (n-3) + 2e^{\gamma/3}$,
since $e^{\gamma/3} > 1$; equality forces $k = 2$, that is $T = P_n$, which
attains the bound with two $(1,2)$-edges and $n-3$ $(2,2)$-edges. For $\gamma < 0$ every
inequality reverses.
\end{proof}

Proposition~\ref{prop:extremal} is a consequence of the criteria of Cruz and Rada~\cite{CruzRada2019PathStar}; the direct proof above is included because it is short. For the weights $\varphi_{ij} = e^{\gamma q_{ij}}$ and $E = e^{\gamma/3}$, their function $f(i,j) = \tfrac{ij}{i+j}(\varphi_{ij} + 2\varphi_{12} - 3\varphi_{22})$ equals $2(E - 1)$ at $(1,2)$ and $(2,2)$. For $\gamma > 0$ it is larger at every other pair: if $i, j \ge 2$ then $\varphi_{ij} + 2E - 3 \ge 2(E-1) > 0$ and $ij/(i+j) > 1$, and for $(1,j)$, $j \ge 3$, one has $f \ge \tfrac34(e^{\gamma/2} + 2E - 3) > 2(E-1)$, since with $x = e^{\gamma/6} > 1$ this reads $3x^3 - 2x^2 - 1 = (x-1)(3x^2 + x + 1) > 0$. Theorem~2.1 of~\cite{CruzRada2019PathStar} then gives the lower bound with equality only for $P_n$, and their Theorem~3.1 gives the upper bound with equality only for $S_n$, because $e^{\gamma q}$ is largest at the pair $(1,n-1)$. For $\gamma < 0$ every inequality reverses. Among chemical trees the bound
$LO(T;0,0,\gamma) \le (n-1)e^{3\gamma/5}$ for $\gamma > 0$ is attained only
for $n = 5$, and the maximal chemical trees for $n > 5$ are left open.

\section{QSPR diagnostics}\label{sec:sensitivity}

We compare the family with the classical indices it reduces to on the
$18$ octane isomers, following the benchmark of the diminished Sombor
study~\cite{MovahediGutmanRedzepovicFurtula2026}, and use the $35$ nonane
isomers as a replication set. No parameter in the correlation comparison
is tuned against the property data. The repository history begins with the octane correlations already
computed, so it cannot document that the benchmark points $\gamma = 1$ and
$2$ were fixed before the octane data were seen. The sets are small, so the octane
correlations are descriptive. After the nonane result of
Section~\ref{sec:ablation} we made the pre-registered tests on measured
boiling points of Section~\ref{sec:bp} the primary predictive evidence.

\subsection{Correlations on the octane isomers}\label{sec:octanes}

Table~\ref{tab:lo_octane} reports Pearson correlations with five
properties. The property values are the standard benchmark compilation,
audited against the NIST WebBook~\cite{NISTWebBook} and the API Research
Project~44 tables~\cite{APIProject44Circ461}. Seven entropies and the
compiled acentric factors rest on secondary sources, and the entropies can
move the $0.001$ margin between $HM$ and $M_1$, so the ordering of the $S$
column is not interpreted. The vaporisation enthalpy of
2,2,3,3-tetramethylbutane, which is crystalline at $298$~K, has no stated
source; with the isomer excluded $LO(0,0,1)$ stays first, level at three
decimals with $LO(0,0,2)$ ($|r| = 0.980$), and with the
hypothetical-liquid estimate it gives $|r| = 0.976$, while the NIST
sublimation enthalpy lowers every $|r|$ to $0.374$ or below, so this column
is conditional on a liquid-phase value. Deleting single isomers shifts each
bolded $|r|$ by at most $0.027$. The descriptors are nearly collinear:
fourteen pairs have $|r| \ge 0.99$, the first two principal components
carry $99.05\%$ of the variance, and $LO(0,0,2)$ has $|r| = 0.993$ with a
classical reduction. Differences in the third decimal are therefore not
rankings, and the bolded entries, each the largest of twelve correlations
on $18$ molecules, are optimistic. The two collision pairs of
Table~\ref{tab:collisions} differ by $1.2\,^{\circ}$C and $2.1\,^{\circ}$C
in $T_B$ although every vertex-degree-based index gives them equal values.

\begin{table}[htbp]
\centering
\caption{Pearson correlation $r$ between the Loyola reductions, the fixed pure-$\gamma$ benchmark points and the five physicochemical properties of the $18$ octane isomers. Bold flags the largest $|r|$ per column, and this is descriptive. Margins to the runner-up are $0.001$--$0.032$, and their paired bootstrap intervals include zero (Supplementary Section~S1). These intervals are post-selection and not adjusted for multiplicity. The entry $-0.000$ ($ISI$, $\Delta H_f$) is $-0.0002$.}
\label{tab:lo_octane}
\small
\begin{tabular}{lrrrrr}
\toprule
Index & $T_B$ & $\Delta H_f$ & $\Delta H_{\mathrm{vap}}$ & $S$ & $\omega$\\
\midrule
$M_1 = LO(0,1,0)$ & $-0.718$ & $-0.762$ & $-0.935$ & $-0.970$ & $-0.970$\\
$M_2 = LO(1,0,0)$ & $-0.497$ & $-0.542$ & $-0.811$ & $-0.934$ & $\mathbf{-0.988}$\\
$HM = LO(0,2,0)$ & $-0.654$ & $-0.710$ & $-0.903$ & $\mathbf{-0.971}$ & $-0.982$\\
${}^{m}\!M_2 = LO(-1,0,0)$ & $\mathbf{+0.855}$ & $\mathbf{+0.891}$ & $+0.925$ & $+0.863$ & $+0.794$\\
$R = LO(-1/2,0,0)$ & $+0.819$ & $+0.850$ & $+0.958$ & $+0.932$ & $+0.898$\\
$\chi = LO(0,-1/2,0)$ & $+0.800$ & $+0.832$ & $+0.961$ & $+0.947$ & $+0.924$\\
$H/2 = LO(0,-1,0)$ & $+0.821$ & $+0.849$ & $+0.961$ & $+0.932$ & $+0.900$\\
$ISI = LO(1,-1,0)$ & $-0.005$ & $-0.000$ & $-0.383$ & $-0.593$ & $-0.738$\\
$GA = 2\,LO(1/2,-1,0)$ & $+0.822$ & $+0.858$ & $+0.965$ & $+0.939$ & $+0.906$\\
$AG = \tfrac12 LO(-1/2,1,0)$ & $-0.816$ & $-0.859$ & $-0.957$ & $-0.938$ & $-0.900$\\
$LO(0,0,1)$ & $-0.828$ & $-0.832$ & $\mathbf{-0.984}$ & $-0.928$ & $-0.929$\\
$LO(0,0,2)$ & $-0.829$ & $-0.846$ & $-0.981$ & $-0.937$ & $-0.925$\\
\bottomrule
\end{tabular}
\end{table}

\subsection{Freeing $\gamma$ and the search box}\label{sec:ablation}

We ask whether, at identical search budgets, freeing $\gamma$ improves
out-of-sample prediction over an optimised parent $M_{\alpha,\beta}$, in a
fully nested leave-one-out design with random search (Supplementary
Section~S3); a parameter-optimised parent has been fitted to octane
properties before~\cite{AliGutmanFurtulaAlbalahiHamza2025}. In the box
$(-2,2)$ freeing $\gamma$ improved the nested $Q^2$ for octane
$\Delta H_{\mathrm{vap}}$ in $100/100$ candidate streams at budgets $200$
and $500$, with median $\Delta Q^2 = +0.069$ and $+0.068$, and in $93/100$
and $96/100$ with 2,2,3,3-tetramethylbutane excluded. It usually degraded
$\Delta H_f$ and was mixed elsewhere. For $\Delta H_{\mathrm{vap}}$, however,
the parent chose $(\alpha,\beta) = (1.021, -1.993)$, on the boundary of the
box, in $16$ of the $18$ folds. By~\eqref{eq:factor} the parent reaches
imbalance information through $(1-q^2)^{\alpha}$ along
$\beta \approx -2\alpha$, which needs large exponents and is the direction
the box truncates. A box of half-width $12$, chosen after inspecting the
parent's nested $Q^2$ curves on the octanes and the nonanes as the
smallest examined in which its selections are interior, was used below for
every property. It does not favour the parent uniformly: for $T_B$ it
lowers the parent's nested $Q^2$ from $0.665$ to $0.606$ on the octanes.

The matched comparison was repeated without pre-registration, with
$2000$ candidates per model and seeds $0$--$49$. In the wide box the
vaporisation-enthalpy gain shrinks to about zero: a hybrid control that
widens only $(\alpha,\beta)$ and keeps $\gamma \in (-2,2)$ gives median
$\Delta Q^2 = +0.004$ on the octanes ($28$ of $50$ seeds) and $-0.002$ on
the nonanes ($24$ of $50$), while widening all three coordinates gives
$-0.010$ and $-0.016$. Most of the change comes from the parent improving
in the wide box ($0.887 \to 0.942$ on the octanes). Octane $S$ moves from
$11$ to $38$ wins, which shows how strongly the win count depends on the
search protocol. Equal budgets do not give equally thorough searches in
two and three dimensions, so the design may cost the three-parameter
model part of the effect sought.

Under a protocol committed before the nonane data were committed and
before any model was fitted to them (commit \texttt{964b127}), the
narrow-box design was applied to the $35$ nonane isomers. The $35$ boiling
points and $34$ standard vaporisation enthalpies are from the NIST
WebBook~\cite{NISTWebBook}; $29$ of the enthalpies are compiled values,
$4$ listed under Labbauf et al.~\cite{LabaufGreenshieldsRossini1961} and
$25$ under Reid~\cite{Reid1972}, a review of the
handbook~\cite{WilhoitZwolinski1971}, so agreement with them measures
consistency with compiled values. The pre-registered verdict is that the
$\gamma$ effect \emph{replicates} for $\Delta H_{\mathrm{vap}}$, with LO
better than GM in $95/100$ and $98/100$ seeds and median
$\Delta Q^2 = +0.048$ and $+0.053$, and is \emph{mixed} for $T_B$ ($58/100$
and $63/100$). We report the verdict as obtained. Its seed-count rule
counts optimiser variance on fixed data, which we regard as a design
error. A post hoc molecule bootstrap conditional on the fitted predictions gives a
$95\%$ interval for LO minus GM on nonane $\Delta H_{\mathrm{vap}}$ of
$[-0.006, +0.113]$ at budget $200$, which includes zero, and in the wide
box, which is not pre-registered, the gain shrinks to about zero. The
untuned point $LO(0,0,1)$, which correlates at $0.996$ with $IRLA$ on both
sets, matches or beats the tuned Loyola model in point estimate ($0.957$ and $0.903$ against
medians $0.953$ and $0.900$; on the octanes this point was singled out
from the same data). Ridge regression on the degree-pair counts has a
higher point estimate than every index model on the nonanes
($Q^2 = 0.927$ for $\Delta H_{\mathrm{vap}}$, $0.682$ for $T_B$; the conditional
interval for LO minus ridge on nonane $\Delta H_{\mathrm{vap}}$, $[-0.101, +0.032]$,
includes zero), but on the
octanes it is below the best single index for four of the five properties
and level on $\Delta H_{\mathrm{vap}}$ ($0.954$ against $0.957$; Supplementary
Table~S7).

\subsection{Measured boiling points}\label{sec:bp}

Two tests used measured boiling points from the NIST
WebBook~\cite{NISTWebBook}, extracted by a rule committed to the
repository (commit \texttt{a5cf378}) before any model was fitted. The rule
admits all hexanes, heptanes and octanes, $34$ of the $35$ nonanes and
$34$ of the $75$ decanes. The included and excluded decanes do not differ
detectably in branching or imbalance (Supplementary Section~S9), but with
$34$ and $41$ molecules this does not show that the decane cohort is
representative. Test~A refits and reselects inside each leave-one-out fold
on the decanes. Test~B pools all $100$ molecules with the carbon number
$n_C$ as a covariate; the octane and nonane boiling points were used
above, so Test~B is a pooled size-adjusted analysis rather than a fresh
external test. Both use the wide boxes, $(\alpha,\beta) \in (-12,12)^2$ and
$\gamma \in (-2,2)$, with $2000$ candidates and seeds $0$--$49$. The
pre-registered criterion for a $\gamma$ effect required at least $40$ wins
in $50$ seeds and a molecule-bootstrap interval excluding zero, which we
read as requiring the lower end of the two-sided $95\%$ interval to exceed
zero.

\begin{table}[htbp]
\centering
\caption{Measured boiling points (nested leave-one-out $Q^2$). Test A:
decane cohort, refitted within folds. Test B: pooled C6--C10 with the carbon number
$n_C$ in every index model (Class., $IRLA$,
LO$_1$: $n_C$ plus that descriptor); Ridge uses the counts alone, whose sum is $n_C - 1$. Base: training-fold mean (Test A)
or $n_C$ alone (Test B).
Class.: best classical index chosen in each fold. LO$_1$: fixed
$LO(0,0,1)$. Ridge: ridge regression on the degree-pair counts. GM,
LO$_h$: medians over seeds $0$--$49$ in the wide boxes, LO$_h$
being the hybrid control with $(\alpha,\beta)$ in the wide box and
$\gamma \in (-2,2)$.
$\Delta Q^2$: median over seeds of the paired difference LO$_h$ minus GM. The registered intervals are conditional on the fitted predictions and omit refitting and selection uncertainty (see below). Last column: seeds (of $50$) in
which LO$_h$ beats GM.}
\label{tab:bp}
\scriptsize
\setlength{\tabcolsep}{1.5pt}
\begin{tabular}{lrrrrrrrrr}
\toprule
Test & Base & Class. & $IRLA$ & LO$_1$ & Ridge & GM & LO$_h$ & $\Delta Q^2$ & $>$\\
\midrule
decanes ($n = 34$) & $-0.062$ & $0.176$ & $0.167$ & $0.149$ & $0.860$ & $0.514$ & $0.406$ & $-0.072$ & $15$\\
C6--C10, with $n_C$ ($n = 100$) & $0.938$ & $0.967$ & $0.963$ & $0.963$ & $0.991$ & $0.977$ & $0.977$ & $-0.0001$ & $25$\\
\bottomrule
\end{tabular}
\end{table}

Both pre-registered verdicts are negative. On the decanes the median
$\Delta Q^2 = -0.072$, LO is ahead in $15$ of $50$ seeds, and the $95\%$
interval is $[-0.156, +0.001]$; on the pooled set the values are
$-0.0001$, $25$ of $50$ and $[-0.003, +0.002]$. In the original box of
half-width $2$, reported as the pre-registration requires, the two tuned
models are level on the decanes (medians $0.469$ and $0.465$), and the
Loyola model is marginally ahead on the pooled set ($0.980$ against
$0.976$); seed-matched comparisons of that narrow-box Loyola model with the
parent were not registered, carry no interval and count optimiser
variance, so they do not indicate a gain from $\gamma$. The registered
intervals are conditional on the fitted predictions and omit refitting and
selection. An exploratory resampling that reruns the whole nested procedure
on random $80\%$ subsamples gives, on the decanes ($500$ subsamples),
median $-0.110$ with central $95\%$ range $[-0.510, +0.312]$, and on the
pooled set ($200$ subsamples) $+0.001$ with range $[-0.015, +0.008]$;
these ranges are not confidence intervals for the full-sample
$\Delta Q^2$. Each replicate uses a single candidate stream, so the ranges
also contain optimiser variance. No smallest effect of interest or equivalence margin was
registered. The negative verdicts therefore mean that no gain was
demonstrated under these protocols; they do not establish equivalence, the
decane range neither supports nor excludes gains of practical size, and
they do not show that $\gamma$ never helps. On the registered secondary
question of Test~B, $IRLA$ and $LO(0,0,1)$ do not beat size plus the best
classical index ($0.963$ against $0.967$), while both tuned models have
higher point estimates ($0.977$; no interval was registered for this
secondary question).

Ridge regression on the ten degree-pair counts reaches $Q^2 = 0.860$ on
the decanes, where every single index stays between $0.149$ and $0.514$
(medians over seeds for the tuned models),
and $0.991$ on the pooled set. A permutation check that repeats the whole nested
procedure, including the choice of the ridge penalty, for $200$
permutations of the decane boiling points never reaches $0.860$ (median
$-0.09$, maximum $0.20$; empirical $p = 1/201 \approx 0.005$). Wiener's
two-descriptor model~\cite{Wiener1947}, with the path number $W$ and the
number $p_3$ of carbon pairs three bonds apart, fitted without tuning,
gives $Q^2 = 0.904$ on the decanes. On trees
$p_3 = \sum_{uv}(d_u - 1)(d_v - 1)$ is a linear function of the degree-pair
counts, so the distance information in this model is $W$. This single
point estimate, given without an interval and for descriptors historically
chosen on alkane boiling points, suggests that within one carbon number
$W$ may add to the degree-pair information; its difference from ridge is
not resolved here.
Under the stated boxes and candidate budgets, ridge regression predicts
the boiling points more accurately than every single-index model
evaluated, with or without $\gamma$, in point estimate; the margin is
large on the decanes and small on the pooled set, where size alone gives
$Q^2 = 0.938$, and no interval for these differences is reported here. These conclusions concern acyclic
alkanes.

\section{Discussion and conclusion}\label{sec:conclusion}

Every index that depends on an edge only through its degree ratio sees a
graph only through its histogram of normalised imbalances. For trees this
histogram carries the full degree-pair information up to order twelve,
and loses it at order thirteen, and at sixteen for chemical trees, with
explicit colliding families. The pure-imbalance points of the Loyola
family at nonzero algebraic $\gamma$ attain this limit exactly; along the
imbalance axis two chemical graphs exchange order at most five times; the
extremal trees of the pure-imbalance points are the path and the star;
and at rational exponents a nonzero algebraic $\gamma$ never destroys a
distinction the parent makes. The minimality of the collision orders,
equality in Proposition~\ref{prop:identrank}(ii) for $\Delta \le 6$ and
the recovery of discrimination off $\Pi$ rest on exhaustive enumeration;
no general argument is given.

On alkanes no predictive gain from $\gamma$ was demonstrated under the
stated protocols, which may handicap the three-parameter search, and on
measured boiling points ridge regression on the degree-pair counts
had higher point estimates than every single index evaluated. The data sets are
small and acyclic: $18$ octanes, $35$ nonanes, $34$ decanes and $100$
pooled C6--C10 alkanes. Seven octane entropies and the compiled acentric
factors rest on secondary sources, and the vaporisation enthalpy of
2,2,3,3-tetramethylbutane has an unresolved reference state. Changes made after outcomes had been seen are
listed in Table~\ref{tab:deviations}; none changes a pre-registered
verdict as reported, but the wide-box analysis that reverses the nonane
conclusion is itself not pre-registered. Open problems are whether five
crossings occur for $\Delta \le 4$, the maximal chemical trees of the
pure-imbalance points, and the behaviour on other compound classes.

\begin{table}[htbp]
\centering
\caption{Changes to the pre-registered protocols made after outcomes
had been seen. Protocols and change records are in the
reproducibility package~\cite{LoyolaRepo}.}
\label{tab:deviations}
\scriptsize
\setlength{\tabcolsep}{3pt}
\begin{tabular}{p{0.2\linewidth}p{0.42\linewidth}p{0.3\linewidth}}
\toprule
Test & Change after outcomes & Effect\\
\midrule
Nonanes (Section~\ref{sec:ablation}) & Parent box widened from half-width
$2$ to $12$, chosen after inspecting the $\Delta H_{\mathrm{vap}}$ curves &
Registered verdict reported as obtained; the gain shrinks to about
zero in the wide box\\
Nonanes & Seed-count criterion judged to count optimiser variance &
None on reported numbers\\
Nonanes & Measured boiling points made the primary predictive evidence &
None on reported numbers\\
Boiling points, Test~B & All NIST pages of a skeleton pooled; one
nonane excluded, one heptane value moved by $0.2\,^{\circ}$C &
Verdict unchanged; both cohorts reported\\
Nonanes & A second NIST $T_B$ of 3-ethyl-4-methylhexane, $6.4$\,K
lower, found later & Still mixed if excluded or replaced (Suppl.~S4)\\
Boiling points, Test~B & Least-squares solver rescaled by columns &
Verdict unchanged\\
Boiling points, Test~B & Scale-invariant rank test & No displayed value changes\\
Boiling points, Test~B & Stricter rank threshold and a deletion-rank guard in the size-adjusted scorer & No output changes\\
Boiling points, both & Interval rule read as one-sided & None\\
Boiling points, both & Selection-aware resampling added
(Section~\ref{sec:bp}) & Exploratory; no verdict depends on it\\
\bottomrule
\end{tabular}
\end{table}

\medskip
\noindent\textbf{Data availability.} Data, scripts and verifiers that
recompute every tabulated value are in the reproducibility
package~\cite{LoyolaRepo}, where the nonane and boiling-point
protocols were committed before the corresponding models were fitted.
The NIST WebBook pages from which the experimental values were
extracted are not redistributed; a script of the package downloads
them and checks each against the copy from which the data sets were
built. Lean~4 with Mathlib checks the convexity of the log-sum-exp form
of $\log LO$, the spanning statement of the independence lemma and the
order-thirteen profile collision on the pure-$\gamma$ line; the
transcendence, Laguerre and enumeration arguments are not
machine-checked.

\acknowledgment{The authors thank colleagues for helpful discussions.}

\singlespacing

 \begin{thebibliography}{00}

    \bibitem{Albertson1997} M. O. Albertson, The irregularity of a graph,
  \textit{Ars Combin.\/} \textbf{46} (1997) 219--225.

  \bibitem{AliGutmanFurtula2026BID} A. Ali, I. Gutman, B. Furtula,
  Extremal results on bond incident degree indices of graphs: a survey,
  \textit{Aequationes Math.\/} \textbf{100} (2026) \#66.

  \bibitem{AliGutmanFurtulaAlbalahiHamza2025} A. Ali, I. Gutman,
  B. Furtula, A. M. Albalahi, A. E. Hamza, On chemical and mathematical
  characteristics of generalized degree-based molecular descriptors,
  \textit{AIMS Math.\/} \textbf{10} (2025) 6788--6804.

  \bibitem{AliGutmanEtAl2027DRSO} A. Ali, I. Gutman, G. G. Nyauli, J. Buragohain,
  F. A. Alhubairah, H. Saber, A. A. Attiya, Degree-ratio Sombor index,
  \textit{MATCH Commun. Math. Comput. Chem.\/} \textbf{97} (2027) 639--656.

  \bibitem{Baker1975} A. Baker, \textit{Transcendental Number Theory\/},
  Cambridge Univ. Press, Cambridge, 1975.

  \bibitem{Besicovitch1940} A. S. Besicovitch, On the linear independence
  of fractional powers of integers, \textit{J. London Math. Soc.\/}
  \textbf{15} (1940) 3--6.

  \bibitem{BollobasErdos1998}
  B. Bollob\'as, P. Erd\H{o}s, Graphs of extremal weights, \textit{Ars Combin.\/}
  \textbf{50} (1998) 225--233.

  \bibitem{CruzMonsalveRada2020} R. Cruz, J. Monsalve, J. Rada, Trees with
  maximum exponential Randi\'c index, \textit{Discrete Appl. Math.\/}
  \textbf{283} (2020) 634--643.

  \bibitem{CruzRada2019PathStar} R. Cruz, J. Rada, The path and the star as
  extremal values of vertex-degree-based topological indices among trees,
  \textit{MATCH Commun. Math. Comput. Chem.\/} \textbf{82} (2019) 715--732.

  \bibitem{DasMondal2024ExpGA} K. C. Das, S. Mondal, On exponential
  geometric-arithmetic index of graphs, \textit{J. Math. Chem.\/}
  \textbf{62} (2024) 2740--2760.

  \bibitem{DasMondalRazaPal2025ExpISDD} P. Das, S. Mondal, Z. Raza,
  A. Pal, On exponential inverse symmetric division deg index of graphs,
  \textit{J. Appl. Math. Comput.\/} \textbf{71} (Suppl.~1) (2025) 537--561.

  \bibitem{DeutschKlavzar2015} E. Deutsch, S. Klav\v{z}ar, M-polynomial and
  degree-based topological indices, \textit{Iran. J. Math. Chem.\/} \textbf{6}
  (2015) 93--102.

  \bibitem{GaoAMC2024} W. Gao, Y. Gao, The extremal trees for
  exponential vertex-degree-based topological indices, \textit{Appl.
  Math. Comput.\/} \textbf{472} (2024) \#128634.

  \bibitem{Gutman2021Sombor} I. Gutman, Geometric approach to degree-based
  topological indices: Sombor indices, \textit{MATCH Commun. Math. Comput.
  Chem.\/} \textbf{86} (2021) 11--16.

  \bibitem{Gutman2026SomborSurvey} I. Gutman, Survey of Sombor-type
  degree-based topological indices, \textit{MATCH Commun. Math. Comput.
  Chem.\/} \textbf{96} (2026) 819--837.

  \bibitem{GutmanBuragohainNyauli2026} I. Gutman, J. Buragohain,
  G. G. Nyauli, Concealed relations between vertex-degree-based
  topological indices, \textit{J. Math. Chem.\/} \textbf{64} (2026) \#39.

  \bibitem{GutmanMilovanovic2020AKCE} I. Gutman, E. Milovanovi\'c,
  I. Milovanovi\'c, Beyond the Zagreb indices, \textit{AKCE Int. J.
  Graphs Comb.\/} \textbf{17} (2020) 74--85.

  \bibitem{GutmanRedzepovicBuragohainNyauli2027} I. Gutman, I. Red\v{z}epovi\'c,
  J. Buragohain, G. G. Nyauli, Comparing degree-ratio Sombor and Sombor indices,
  \textit{MATCH Commun. Math. Comput. Chem.\/} \textbf{97} (2027) 165--174.

  \bibitem{GutmanRedzepovicKizilirmakKulli2025EU} I. Gutman, I. Red\v{z}epovi\'c,
  G. \"Ozkan Kizilirmak, V. R. Kulli, Euler--Sombor index and its congeners,
  \textit{Open J. Math. Sci.\/} \textbf{9} (2025) 141--148.

  \bibitem{GutmanTrinajstic1972} I. Gutman, N. Trinajsti\'c, Graph theory
  and molecular orbitals. Total $\pi$-electron energy of alternant
  hydrocarbons, \textit{Chem. Phys. Lett.\/} \textbf{17} (1972) 535--538.

  \bibitem{HuMATCH2022} Z. Hu, L. Li, X. Li, D. Peng, Extremal graphs
  for topological index defined by a degree-based edge-weight
  function, \textit{MATCH Commun. Math. Comput. Chem.\/} \textbf{88}
  (2022) 505--520.

  \bibitem{Jameson2006} G. J. O. Jameson, Counting zeros of generalised
  polynomials: Descartes' rule of signs and Laguerre's extensions,
  \textit{Math. Gaz.\/} \textbf{90} (2006) 223--234.

  \bibitem{Konstantinova1996Degeneracy} E. V. Konstantinova, The
  discrimination ability of some topological and information distance indices
  for graphs of unbranched hexagonal systems, \textit{J. Chem. Inf. Comput.
  Sci.\/} \textbf{36} (1996) 54--57.

  \bibitem{LabaufGreenshieldsRossini1961} A. Labbauf, J. B. Greenshields,
  F. D. Rossini, Heats of formation, combustion, and vaporization of the
  35 nonanes and 75 decanes, \textit{J. Chem. Eng. Data\/} \textbf{6}
  (1961) 261--263.

  \bibitem{LinZhouWangMiao2022} Z. Lin, T. Zhou, X. Wang, L. Miao, The general
  Albertson irregularity index of graphs, \textit{AIMS Math.\/} \textbf{7}
  (2022) 25--38.

  \bibitem{NISTWebBook} P. J. Linstrom, W. G. Mallard (Eds.), \textit{NIST
  Chemistry WebBook, NIST Standard Reference Database Number 69\/}, National
  Institute of Standards and Technology, Gaithersburg, 2025,
  \url{https://doi.org/10.18434/T4D303}.

  \bibitem{Mityagin2020} B. S. Mityagin, The zero set of a real analytic
  function, \textit{Math. Notes\/} \textbf{107} (2020) 529--530.

  \bibitem{MolinaSigarreta2025GM} E. D. Molina, J. M.
  Rodr\'iguez-Garc\'ia, J. M. Sigarreta, S. J. Torralbas Fitz, On the
  Gutman--Milovanovi\'c index and chemical applications, \textit{AIMS
  Math.\/} \textbf{10} (2025) 1998--2020.

  \bibitem{MovahediGutmanRedzepovicFurtula2026} F. Movahedi, I. Gutman, I.
  Red\v{z}epovi\'c, B. Furtula, Diminished Sombor index, \textit{MATCH Commun.
  Math. Comput. Chem.\/} \textbf{95} (2026) 141--162.

  \bibitem{LoyolaRepo} A. Natarajan, G. Shiwakoti, M. Arockiaraj,
  The Loyola index family: data, scripts and verifiers,
  https://github.com/Singati2/LOYOLA-PAPER (2026).

  \bibitem{NyauliBuragohain2026CoRSO} G. G. Nyauli, J. Buragohain, Geometric
  approach to degree-based topological indices: cosine-rule generalized
  Sombor indices, \textit{MATCH Commun. Math. Comput. Chem.\/} \textbf{96}
  (2026) 295--324.

  \bibitem{NyauliBuragohain2027DRSO} G. G. Nyauli, J. Buragohain,
  Geometric approach to degree-based topological indices: degree-ratio
  Sombor indices of graphs, \textit{MATCH Commun. Math. Comput. Chem.\/}
  \textbf{97} (2027) 135--164.

  \bibitem{Rada2019ExpVDB} J. Rada, Exponential vertex-degree-based
  topological indices and discrimination, \textit{MATCH Commun. Math.
  Comput. Chem.\/} \textbf{82} (2019) 29--41.

  \bibitem{Rada2026LinearVDB} J. Rada, A finite-dimensional linear
  framework for vertex-degree-based topological indices, \textit{MATCH
  Commun. Math. Comput. Chem.\/} \textbf{96} (2026) 841--863.

  \bibitem{RadaRodriguezSigarreta2022Weak} J. Rada, J. M. Rodr\'{\i}guez,
  J. M. Sigarreta, Weakly discriminating vertex-degree-based topological
  indices, \textit{MATCH Commun. Math. Comput. Chem.\/} \textbf{87} (2022)
  97--108.

  \bibitem{Randic1975} M. Randi\'c, Characterization of molecular
  branching, \textit{J. Am. Chem. Soc.\/} \textbf{97} (1975) 6609--6615.

  \bibitem{Reid1972} R. C. Reid, Review of R. C. Wilhoit, B. J. Zwolinski, Handbook
  of Vapor Pressures and Heats of Vaporization of Hydrocarbons and Related
  Compounds, \textit{AIChE J.\/} \textbf{18} (1972) 1278.

  \bibitem{RetiSharafdini2018IRLA} T. R\'eti, R. Sharafdini,
  \'A. Dr\'egelyi-Kiss, H. Haghbin, Graph irregularity indices used as
  molecular descriptors in QSPR studies, \textit{MATCH Commun. Math.
  Comput. Chem.\/} \textbf{79} (2018) 509--524.

  \bibitem{APIProject44Circ461} F. D. Rossini, K. S. Pitzer, W. J. Taylor,
  J. P. Ebert, J. E. Kilpatrick, C. W. Beckett, M. G. Williams, H. G. Werner,
  \textit{Selected Values of Properties of Hydrocarbons\/}, Natl. Bur.
  Stand. Circ. 461, U.S. Government Printing Office, Washington, 1947.

  \bibitem{ShaoJahanbani2027} Z. Shao, A. Jahanbani, On the degree-ratio Sombor
  index: corrections and new bounds, \textit{MATCH Commun. Math. Comput.
  Chem.\/} \textbf{97} (2027) 657--675.

  \bibitem{ShirdelRezapourSayadi2013HM} G. H. Shirdel, H. Rezapour,
  A. M. Sayadi, The hyper-Zagreb index of graph operations,
  \textit{Iran. J. Math. Chem.\/} \textbf{4} (2013) 213--220.

  \bibitem{Sigarreta2022ExpVDB} J. M. Sigarreta, Extremal problems on
  exponential vertex-degree-based topological indices, \textit{Math.
  Biosci. Eng.\/} \textbf{19} (2022) 6985--6995.

  \bibitem{SongHuang2023EDV} R. Song, Q. Huang, A relation on trees and the
  topological indices based on subgraph, \textit{MATCH Commun. Math.
  Comput. Chem.\/} \textbf{89} (2023) 343--370.

  \bibitem{VukicevicFurtula2009GA} D. Vuki\v{c}evi\'c, B. Furtula, Topological
  index based on the ratios of geometrical and arithmetical means of
  end-vertex degrees of edges, \textit{J. Math. Chem.\/} \textbf{46} (2009)
  1369--1376.

  \bibitem{VukicevicGasperov2010} D. Vuki\v{c}evi\'c, M. Ga\v{s}perov, Bond additive
  modeling 1. Adriatic indices, \textit{Croat. Chem. Acta\/} \textbf{83} (2010)
  243--260.

  \bibitem{Wiener1947} H. Wiener, Structural determination of paraffin boiling
  points, \textit{J. Am. Chem. Soc.\/} \textbf{69} (1947) 17--20.

  \bibitem{WilhoitZwolinski1971} R. C. Wilhoit, B. J. Zwolinski,
  \textit{Handbook of Vapor Pressures and Heats of Vaporization of
  Hydrocarbons and Related Compounds\/}, Thermodynamics Research Center, Texas A\&M Research Foundation, College Station, 1971.

  \bibitem{ZhouTrinajstic2009sumconn}
  B. Zhou, N. Trinajsti\'c, On a novel connectivity index, \textit{J. Math.
  Chem.\/} \textbf{46} (2009) 1252--1270.

  \bibitem{ZhouTrinajstic2010} B. Zhou, N. Trinajsti\'c, On general
  sum-connectivity index, \textit{J. Math. Chem.\/} \textbf{47} (2010)
  210--218.

\end{thebibliography}

\end{document}

